\section{结论}
\label{sec:conclusion}

本文指出现代 MLLM 的一个结构性问题：\textbf{跨模态粒度失配}。连续视觉特征与离散、碎片化的子词 token 之间的不一致，使自回归生成易处于未经验证的\emph{开环解码}状态，语言先验惯性常压倒视觉约束，从而催化物体与关系幻觉。

针对该根因，我们提出 \textbf{CHORD}（Calibrating Hallucinations via Object-Resonant Decoding）：通过无参数分词器桥接将碎片 token 提升为语义完整片段，在外部共享语义空间中计算\textbf{共振增量 $\Delta S$}，在避免脆弱启发式与沉重黑盒评估器的同时，强化视觉一致内容并抑制虚构描述。

我们进一步用配备 KV 缓存加速的\textbf{连续 Top-$1$ 监控}替代朴素熵门控，在控制开销的同时仍能有效应对过度自信型幻觉。

多宿主家族上的实验表明，CHORD 在显著抑制物体幻觉的同时，尽量保持语言流畅性。结果表明，通过外部物体共振解码可实现强且可跨模型的接地控制，为更可信的 MLLM 部署提供可扩展范式。
